% ===== UNIT 17: Score-level fusion: why score- not feature-level (3.5 part 2) =====
% Status: AI-DRAFT  (change to REWRITTEN / REVIEWED when done; see CHECKLIST.md)
% Rules: style rewrite only. Keep every \cite, \ref, \label, number, equation and table value as is. Change the wording, not the content.
Score-level fusion is preferred over feature-level fusion because the two
streams have fundamentally different feature geometries (1536D concat-patch ViT
representation vs.\ 2048D convolutional map) and normalisation schemes;
combining at the feature level would require a learned projection prone to
overfitting on 576 training identities. The streams also use different input
resolutions (224\,$\times$\,224 vs.\ 320\,$\times$\,320, each matching its
checkpoint); score-level fusion is agnostic to this mismatch by construction,
and harmonising would actively hurt---evaluating Stream~1 at
320\,$\times$\,320 costs $1.7$\,pp mAP (Section~\ref{sec:res_sens}).

We also verified the score-versus-feature-level choice empirically. Under the
identical downstream post-processing, plain L2-normalised concatenation of the
two streams reaches 93.08\% mAP / 98.63\% Rank-1, and \emph{weight-matched}
concatenation ($[\sqrt{w_1}\,\mathbf{f}_1;\sqrt{w_2}\,\mathbf{f}_2]$, the
feature-space counterpart of Eq.~\eqref{eq:fusion}) reaches 93.37\% /
98.51\%---statistically indistinguishable from the deployed score fusion
(93.32\% / 98.45\%)---whereas PCA-whitening (1024-D) fitted strictly on the
training split degrades sharply to 86.50\%. Feature- and score-level fusion are
therefore equivalent here when equally weighted, with no evidence that
feature-level concatenation overfits; we retain score-level fusion for its
interpretability (explicit stream weights) and because classical
whitening-style feature calibration actively hurts in this regime. The optimal
weights $w_1{=}0.35$, $w_2{=}0.65$ were found by grid search on the VeRi-776
query/gallery split ($\approx$49.5\,min on T4). Since no separate held-out
validation split is used for weight selection, the chosen weights may reflect a
mild optimistic bias. However, Figure~\ref{fig:weight} demonstrates that mAP
remains above 93.0\% across $w_1 \in [0.2,\,0.5]$, indicating that the optimal
weight is not a sharp peak and that the practical impact of this selection
procedure on the reported result is negligible.
